%Version 3.1 December 2024
% See section 11 of the User Manual for version history
%
%%%%%%%%%%%%%%%%%%%%%%%%%%%%%%%%%%%%%%%%%%%%%%%%%%%%%%%%%%%%%%%%%%%%%%
%%                                                                 %%
%% Please do not use \input{...} to include other tex files.       %%
%% Submit your LaTeX manuscript as one .tex document.              %%
%%                                                                 %%
%% All additional figures and files should be attached             %%
%% separately and not embedded in the \TeX\ document itself.       %%
%%                                                                 %%
%%%%%%%%%%%%%%%%%%%%%%%%%%%%%%%%%%%%%%%%%%%%%%%%%%%%%%%%%%%%%%%%%%%%%

%%\documentclass[referee,sn-basic]{sn-jnl}% referee option is meant for double line spacing

%%=======================================================%%
%% to print line numbers in the margin use lineno option %%
%%=======================================================%%

%%\documentclass[lineno,pdflatex,sn-basic]{sn-jnl}% Basic Springer Nature Reference Style/Chemistry Reference Style

%%=========================================================================================%%
%% the documentclass is set to pdflatex as default. You can delete it if not appropriate.  %%
%%=========================================================================================%%

%%\documentclass[sn-basic]{sn-jnl}% Basic Springer Nature Reference Style/Chemistry Reference Style

%%Note: the following reference styles support Namedate and Numbered referencing. By default the style follows the most common style. To switch between the options you can add or remove “Numbered” in the optional parenthesis. 
%%The option is available for: sn-basic.bst, sn-chicago.bst%  
 
%%\documentclass[pdflatex,sn-nature]{sn-jnl}% Style for submissions to Nature Portfolio journals
%%\documentclass[pdflatex,sn-basic]{sn-jnl}% Basic Springer Nature Reference Style/Chemistry Reference Style
\documentclass[pdflatex,sn-mathphys-num, oneside, 11pt]{sn-jnl}% Math and Physical Sciences Numbered Reference Style
%%\documentclass[pdflatex,sn-mathphys-ay]{sn-jnl}% Math and Physical Sciences Author Year Reference Style
%%\documentclass[pdflatex,sn-aps]{sn-jnl}% American Physical Society (APS) Reference Style
%%\documentclass[pdflatex,sn-vancouver-num]{sn-jnl}% Vancouver Numbered Reference Style
%%\documentclass[pdflatex,sn-vancouver-ay]{sn-jnl}% Vancouver Author Year Reference Style
%%\documentclass[pdflatex,sn-apa]{sn-jnl}% APA Reference Style
%%\documentclass[pdflatex,sn-chicago]{sn-jnl}% Chicago-based Humanities Reference Style

%%%% Standard Packages
%%<additional latex packages if required can be included here>

\usepackage{graphicx}%
\usepackage{multirow}%
\usepackage{amsmath,amssymb,amsfonts}%
\usepackage{amsthm}%
\usepackage{mathrsfs}%
\usepackage[title]{appendix}%
\usepackage{xcolor}%
\usepackage{textcomp}%
\usepackage{manyfoot}%
\usepackage{booktabs}%
\usepackage{algorithm}%
\usepackage{algorithmicx}%
\usepackage{algpseudocode}%
\usepackage{listings}%
\usepackage{geometry}
\geometry{
  reset,          
  a4paper,
  textwidth=190mm,% Forces the text width (A4 is 210mm wide - 10L - 10R = 190mm)
  textheight=267mm, % Forces the text height (A4 is 297mm tall - 15T - 15B = 267mm)
  left=10mm,
  right=10mm,
  top=15mm,
  bottom=15mm
}
%%%%

%%%%%=============================================================================%%%%
%%%%  Remarks: This template is provided to aid authors with the preparation
%%%%  of original research articles intended for submission to journals published 
%%%%  by Springer Nature. The guidance has been prepared in partnership with 
%%%%  production teams to conform to Springer Nature technical requirements. 
%%%%  Editorial and presentation requirements differ among journal portfolios and 
%%%%  research disciplines. You may find sections in this template are irrelevant 
%%%%  to your work and are empowered to omit any such section if allowed by the 
%%%%  journal you intend to submit to. The submission guidelines and policies 
%%%%  of the journal take precedence. A detailed User Manual is available in the 
%%%%  template package for technical guidance.
%%%%%=============================================================================%%%%

%% as per the requirement new theorem styles can be included as shown below
\theoremstyle{thmstyleone}%
\newtheorem{theorem}{Theorem}%  meant for continuous numbers
%%\newtheorem{theorem}{Theorem}[section]% meant for sectionwise numbers
%% optional argument [theorem] produces theorem numbering sequence instead of independent numbers for Proposition
\newtheorem{proposition}[theorem]{Proposition}% 
%%\newtheorem{proposition}{Proposition}% to get separate numbers for theorem and proposition etc.

\theoremstyle{thmstyletwo}%
\newtheorem{example}{Example}%
\newtheorem{remark}{Remark}%

\theoremstyle{thmstylethree}%
\newtheorem{definition}{Definition}%

\raggedbottom
%%\unnumbered% uncomment this for unnumbered level heads

% --- NEW SPACE-SAVING COMMANDS ---
\usepackage{microtype}      % 1. Optimizes horizontal text packing (Saves ~5% space)
\linespread{0.96}           % 2. Reduces vertical space between lines (Saves ~5% space)
\usepackage{newtxtext}      % 3. Ensures you are using the most modern, compact Times font
\usepackage{newtxmath}
% ---------------------------------

\begin{document}

\title[Article Title]{Are meteorological, agricultural, and satellite features from the CY-Bench Dataset good predictors of European Wheat Futures Price Changes?}

%%=============================================================%%
%% GivenName	-> \fnm{Joergen W.}
%% Particle	-> \spfx{van der} -> surname prefix
%% FamilyName	-> \sur{Ploeg}
%% Suffix	-> \sfx{IV}
%% \author*[1,2]{\fnm{Joergen W.} \spfx{van der} \sur{Ploeg} 
%%  \sfx{IV}}\email{iauthor@gmail.com}
%%=============================================================%%

\author{\fnm{68729}}
\author{\fnm{candidate number}}
\author{\fnm{candidate number}}

%%\pacs[JEL Classification]{D8, H51}

%%\pacs[MSC Classification]{35A01, 65L10, 65L12, 65L20, 65L70}

\maketitle

\clearpage

\section{Executive Summary}\label{sec1}

\clearpage

\section{Table of Contents}\label{sec2}

\clearpage

\section{Introduction (1 page)}\label{sec3}

\clearpage

\section{(Optional) Literature Review (1 page)}\label{sec3}

There are two sets of research topics that motivate our literature review: the application of machine learning to multi-national crop yield forecasting, and the use of data-driven models to forecast wheat futures using agricultural signals as inputs. Both areas of research have developed substantially in recent years, but largely in parallel, with limited work connecting yield forecasting outputs to commodity price models. The review below covers both sets of papers.
 
We start with the foundations of ML-based crop yield prediction as described by van Klompenburg, Kassahun, and Catal~\cite{bib14}, who conducted a systematic literature review of 50 studies across a range of crops, algorithms, and geographies. The feature sets and model choices most commonly adopted were temperature, rainfall, and soil type as the dominant predictors. In terms of model selection, their review found that no single model emerged as definitively best, but that Random Forest, Neural Networks, Linear Regression, and Gradient Boosting were the most frequently used, with most studies testing a variety of models to identify the best performer for their specific task. The most widely used evaluation method was 10-fold cross-validation, and artificial neural networks were the most widely applied model class. They also looked at deep learning literature and found CNNs, LSTMs, and DNNs to be the core architectures for this prediction task. Importantly, their review highlighted data availability as the primary constraint on generalization. They found that studies were typically operated over limited geographic extents with short time series, making it difficult to draw conclusions about which methods transfer across regions or climates. Oikonomidis, Catal, and Kassahun~\cite{bib15} extended the ML crop yield review specifically to deep learning approaches, systematically analysing 44 primary studies. Their results confirmed CNN as both the most frequently applied deep learning algorithm and the best-performing one across studies, with LSTM and DNN following as the next most common architectures. They found that RMSE was the dominant evaluation metric used across the reviewed literature. Crucially, they also found that the most significant challenge was the limitation of large training datasets due to insufficient data availability. This raised the risk of overfitting and reduced real-world model performance, with proposed solutions including data augmentation and transfer learning. This data constraint is consistent with what van Klompenburg et al.~\cite{bib14} identified in their ML literature as the primary constraint, so together the two reviews establish data availability and geographic coverage as the primary bottleneck to building generalisable crop yield prediction models. Muruganantham et al.~\cite{bib20} conducted a further systematic review specifically combining deep learning with remote sensing for crop yield prediction, analysing 44 studies from 2012 to 2022 and confirming the previous findings that LSTM and CNN are the dominant approaches, with satellite-derived vegetation indices emerging as among the most informative input features alongside weather variables.
 
Paudel et al.~\cite{bib16} attempt to address these gaps by establishing a large-scale ML baseline for crop yield prediction across multiple countries and crop types. Working with a feature set covering meteorological variables, crop model outputs, and remote sensing indicators, they evaluated Ridge Regression, KNN, SVM, and Gradient Boosting under a time-based look-forward cross-validation scheme that simulates operational forecasting conditions by training only on data preceding each test year on 70 percent train, 30 percent test split. They looked at three European countries, the Netherlands, Germany, and France, evaluating five crop types, with predictions made at the regional level and then aggregated to the national level for comparison with the established European Commission’s "Monitoring Agricultural ResourceS Crop Yield Forecasting System" (MCYFS). Their results showed that early-season predictions made 30 days after planting were broadly competitive with the MCYFS in most cases. They note that accuracy improved as the season progressed toward harvest, confirming that lead time is a key constraint for any real-time application.
 
The most recent CY-Bench paper Paudel et al.~\cite{bib17} built on this foundation at a global scale, combining predictor data, including meteorological variables such as temperature, precipitation, evapotranspiration, and solar radiation, with soil data like available water capacity, bulk density, drainage class, and moisture content. Lastly, they also used satellite vegetation across 29 wheat-growing and 42 maize-growing countries. For features, the dataset includes threshold-based indicators such as counts of cold days, hot days, and dry days during the growing season, for the purpose of tracking climate-related stress signals on the crops. Model evaluation uses Leave One Year Out cross-validation to obtain estimates of performance across both typical and extreme harvest years, with normalized root mean squared error and mean absolute percentage error as the primary metrics to support cross-country comparison. Baseline models evaluated include Ridge Regression, Random Forest, and LSTM, which were all compared against a Naive Average Yield benchmark. These two Paudel studies together define the current benchmark results for the CY-Bench dataset for crop yield forecasting.
 
In terms of modelling agricultural commodity price dynamics, Theofilou et al.~\cite{bib21} provide a systematic review of 50 Machine Learning studies applied to wheat, corn, and rice price forecasting, finding that hybrid deep learning models consistently outperform traditional approaches (such as ARIMA and Partial Least Squares Regression), and identifying persistent challenges including reliance on region-specific datasets and limited model interpretability. Li~\cite{bib18} benchmarks Holt-Winters exponential smoothing against three LSTM-based variants on a European wheat spot price series, with the CNN-LSTM model achieving the best performance due to its ability to capture long-term dependencies and price trend characteristics. Zhang and Tang~\cite{bib22} propose a hybrid LSTM model where they decompose the agricultural commodity futures price time series into multiple smoother components (also called variational mode decomposition), finding that decomposition-based preprocessing substantially improves accuracy for volatile price series. Brignoli et al.~\cite{bib23} compare LSTM-RNNs against classical econometric models for grain futures price forecasting, finding that LSTM-RNNs consistently outperform at longer forecast horizons and are particularly capable of handling structural breaks in price series. They note specific challenges in handling seasonality and trends requiring extra pre-processing steps to remove these factors. Thaker, Chan, and Sonner~\cite{bib19} take a fundamentally different approach that bridges the yield and price literatures directly since they apply a CNN to aerial imagery of winter hard red wheat planted areas. They use cloud coverage patterns as a weather proxy for field-level yield conditions, and then train the model on this imagery to produce a 20 trading-day forecast for wheat futures price on the Chicago Mercantile Exchange with purely directional long or short trading recommendations. Their paper demonstrates that satellite-derived signals can generate positive alpha in wheat futures, but they also document performance decay once the signal becomes widely available. This is, again, a natural and expected limitation of this type of prediction task, potentially showing that the informational advantage of agricultural based price forecasts depend on signal differentiation from other traders.
 
Together, the reviewed literature establishes a clear gap. Indeed, the yield prediction task has produced well-validated, multi-national, in-season yield forecasts but has not connected these directly to commodity futures price modelling. Secondly, the literature on the price modelling task demonstrates the growing effectiveness of deep learning and hybrid models for agricultural price forecasting, but these approaches tend to focus on price history and market signals rather than on benchmarked yield predictors. Finally, Thaker et al.~\cite{bib19} come closest to bridging this gap by using satellite-derived yield proxies as price signals, however, they mainly rely on unvalidated aerial imagery rather than more established predictors. Our study attempts to address this gap by using CY-Bench yield predictors as structured fundamental inputs to a wheat futures price model.

\clearpage

\section{Data Pre-processing (2 pages)}\label{sec3}

\clearpage

\section{Redefining the Prediction Task (1 page)}\label{sec3}

\clearpage

\section{Feature Engineering and Selection (2 pages)}\label{sec3}

\clearpage

\section{Method 1 (2 pages)}\label{sec3}

\clearpage

\section{Method 2 (2 pages)}\label{sec3}

\clearpage

\section{Method 3 (2 pages)}\label{sec3}

\clearpage

\section{Method 4 (2 pages)}\label{sec3}

\clearpage

\section{Model Result Comparison (1 page)}\label{sec3}

\clearpage

\section{Conclusion (1 page)}\label{sec3}

\clearpage

\bibliography{sn-bibliography}% common bib file
%% if required, the content of .bbl file can be included here once bbl is generated
%%\input sn-article.bbl

\clearpage

\section{Generative AI Statement}

\clearpage

\begin{appendices}

\section{Section title of first appendix}\label{secA1}

An appendix contains supplementary information that is not an essential part of the text itself but which may be helpful in providing a more comprehensive understanding of the research problem or it is information that is too cumbersome to be included in the body of the paper.

%%=============================================%%
%% For submissions to Nature Portfolio Journals %%
%% please use the heading ``Extended Data''.   %%
%%=============================================%%

%%=============================================================%%
%% Sample for another appendix section			       %%
%%=============================================================%%

%% \section{Example of another appendix section}\label{secA2}%
%% Appendices may be used for helpful, supporting or essential material that would otherwise 
%% clutter, break up or be distracting to the text. Appendices can consist of sections, figures, 
%% tables and equations etc.

\end{appendices}

%%===========================================================================================%%
%% If you are submitting to one of the Nature Portfolio journals, using the eJP submission   %%
%% system, please include the references within the manuscript file itself. You may do this  %%
%% by copying the reference list from your .bbl file, paste it into the main manuscript .tex %%
%% file, and delete the associated \verb+\bibliography+ commands.                            %%
%%===========================================================================================%%

\end{document}
